\documentclass[10pt,a4paper]{article}
\usepackage[utf8]{inputenc}
\usepackage[spanish]{babel}
\usepackage{url}
\usepackage{hyperref}
\usepackage{amsmath}
\usepackage{amsfonts}
\usepackage{amssymb}
\usepackage{graphicx}
\usepackage{float}
\usepackage{lipsum}
\usepackage{multicol}
\usepackage{xcolor}
\newcommand{\celda}[1]{
	\begin{minipage}{2.5cm}
		\vspace{5mm}
		#1
		\vspace{5mm}
	\end{minipage}
}
\definecolor{azul}{rgb}{0.0, 0.53, 0.74}
\usepackage[left=2.00cm, right=2.00cm, top=2.00cm, bottom=2.00cm]{geometry}
\author{Juan Felipe Calixto Londoño}
\title{Análisis Analítico de Propuestas de Valor en Posgrados de Gerencia y Gestión de Proyectos}

\begin{document}
	
	\begin{figure}[H]
		\raggedright
		% \includegraphics[scale=0.1]{sm.png} \hfill \includegraphics[scale=0.35]{villa.png}
	\end{figure}

	\vspace{7.5mm}
	
	\begin{center}
		{\Large \textbf{Análisis de Propuestas de Valor en Posgrados de Gerencia y Gestión de Proyectos}}\\
		\vspace{2mm}
		{\large Juan Felipe Calixto Londoño$^{1}$}\\
		\vspace{7.5mm}
		$^1$\textit{Universidad Nacional de Colombia. Facultad de ingeniería.} \\
	\end{center}

	\begin{center}
		\textcolor{azul}{\rule{150mm}{0.5mm}}
	\end{center}		

	\begin{abstract}
		El presente documento analiza y compara las propuestas de valor ofrecidas por diversas instituciones de educación superior, tanto a nivel nacional (Colombia) como internacional, en sus programas de posgrado (especializaciones y maestrías) enfocados en la Gerencia y Gestión de Proyectos. A través de este análisis, se identifican las tendencias actuales, tales como la incorporación de metodologías ágiles, la sostenibilidad y el liderazgo, que diferencian a los programas en un mercado altamente competitivo e impulsado por estándares globales como los del Project Management Institute (PMI).
		\vspace{2mm}\\
		\underline{\textbf{Palabras clave:}} \hspace{1mm} \textit{Gerencia de Proyectos, Propuesta de Valor, Posgrados, Educación Superior, PMI, Metodologías Ágiles.}
	\end{abstract}

	\begin{center}
		\textcolor{azul}{\rule{150mm}{0.5mm}}
	\end{center}		

	\vspace{5mm}
	
	\begin{multicols}{2}
		\section{Introducción}
			La Gerencia y Gestión de Proyectos se ha convertido en una disciplina fundamental para el desarrollo estratégico de las organizaciones modernas. Ante esta creciente demanda, las instituciones educativas ofrecen diversos programas de posgrado. Este taller explora cómo estas instituciones estructuran su ``propuesta de valor'' para atraer profesionales, diferenciándose mediante acreditaciones internacionales, enfoques en innovación y fortalecimiento de habilidades blandas (soft skills).
			
		\section{Universidades en Colombia}
			En el contexto colombiano, las propuestas de valor varían según la filosofía de cada institución:
			\begin{itemize}
			    \item \textbf{Universidad de los Andes:} Subraya la formación de líderes estratégicos capaces de alinear proyectos con la visión corporativa y la innovación.
			    \item \textbf{Universidad EAN:} Su diferenciador principal es el enfoque en la sostenibilidad y el emprendimiento, preparando gerentes para proyectos con triple impacto.
			    \item \textbf{Universidad Nacional:} Se enfoca en el rigor académico e investigativo, aportando fuertemente al sector público y de infraestructura civil.
			\end{itemize}
	\end{multicols}
	
	\begin{table}[H]
		\centering
		\begin{tabular}{p{4cm} p{6cm} p{4cm}}
			\hline
			\vspace{2mm}\textbf{Institución}\vspace{2mm} & \vspace{2mm}\textbf{Enfoque Principal}\vspace{2mm} & \vspace{2mm}\textbf{Diferenciador}\vspace{2mm}\\
			\hline
			\vspace{2mm}Univ. de los Andes\vspace{2mm} & \vspace{2mm}Estrategia e Innovación\vspace{2mm} & \vspace{2mm}Alineación PMI\vspace{2mm}\\
			\hline
			\vspace{2mm}Universidad EAN\vspace{2mm} & \vspace{2mm}Sostenibilidad\vspace{2mm} & \vspace{2mm}Proyectos de Triple Impacto\vspace{2mm}\\
			\hline
			\vspace{2mm}Stanford (EE.UU.)\vspace{2mm} & \vspace{2mm}Agilidad e Innovación Continua\vspace{2mm} & \vspace{2mm}Design Thinking\vspace{2mm}\\
			\hline
		\end{tabular}
		\caption{Tabla Comparativa de Propuestas de Valor}
		\label{tb: Tabla1}
	\end{table}

	\begin{multicols}{2}
		\section{Universidades Internacionales}
			Instituciones de renombre mundial establecen las tendencias en gestión de proyectos:
			\begin{itemize}
			    \item \textbf{Stanford University:} Su Advanced Project Management certificate enfatiza metodologías ágiles en entornos de alta incertidumbre y Silicon Valley.
			    \item \textbf{Massachusetts Institute of Technology (MIT):} Integra ingeniería de sistemas complejos y ciencia de datos aplicados a la dirección de proyectos.
			\end{itemize}
		
		\section{Análisis Comparativo}
			Al comparar los currículos, las instituciones colombianas prestan especial atención a la certificación PMP (Project Management Professional), dado el gran valor que tiene en el mercado laboral latinoamericano. Por otro lado, las universidades top internacionales combinan los marcos de trabajo globales (Scrum, Kanban, IPMA, PMI) con tecnologías emergentes e inteligencia artificial aplicadas a la estimación y riesgos.
			
		\section{Conclusiones}
			Las propuestas de valor de los programas de maestría y especialización no se centran hoy en día puramente en cronogramas y presupuestos. El verdadero valor se inclina hacia la creación de líderes adaptables al cambio, la adopción de enfoques híbridos/ágiles y la sostenibilidad.
			
	    \begin{thebibliography}{00}
	    
	    \bibitem{Andes2023}
	        \newblock \textit{Maestría en Gerencia de Proyectos}
	        \newblock Universidad de los Andes.
	        \newblock 2023.
	    
	    \bibitem{PMI2021}
	        \newblock \textit{A Guide to the Project Management Body of Knowledge (PMBOK Guide) – Seventh Edition}
	        \newblock Project Management Institute.
	        \newblock 2021.
	        
	    \end{thebibliography}
	\end{multicols}
	
\end{document}